\documentclass[10pt,a4paper]{amsart}
\usepackage[margin=19mm]{geometry}
\usepackage{amsmath,amssymb,mathtools}
\usepackage[scr=rsfs,cal=euler]{mathalfa}
\usepackage{xfrac}
\usepackage[unicode,hidelinks,bookmarks=false]{hyperref}
\newcommand{\ie}{i.e.\ }
\newcommand{\cf}{cf.\ }
\newcommand{\resp}{respectively\ }
\newcommand{\Subsection}{\subsection}
\providecommand{\nobreakdash}{\nobreak\hbox{-}}
\setlength{\emergencystretch}{2em}
\multlinegap0pt
\usepackage{bm}
\usepackage{bbm}
\usepackage{dsfont}
\usepackage{xspace}
\usepackage{cancel}
\usepackage{mathtools}
\usepackage{amsthm}
\makeatletter
\@ifundefined{thm}{\newtheorem{thm}{Theorem}}{}
\@ifundefined{cor}{\newtheorem{cor}[thm]{Corollary}}{}
\@ifundefined{lem}{\newtheorem{lem}[thm]{Lemma}}{}
\@ifundefined{prop}{\newtheorem{prop}[thm]{Proposition}}{}
\makeatother
\newcommand{\diag}{\textsubscript{\raisebox{0.4ex}{\scalebox{0.6}{\(\blacktriangle\)}}}}\xspace
\title[Locally integrable cross sections and their intersection cov]{Locally integrable cross sections and their intersection covolume}
\author{Nachi Avraham-Re'em, Michael Bj\"orklund,, Rickard Cullman}
\date{Source: arXiv:2509.20836v3 (CC BY 4.0)}
\begin{document}
\maketitle

\noindent\emph{Source and licence.} Locally integrable cross sections and their intersection covolume. Version arXiv:2509.20836v3 (CC BY 4.0), distributed under the Creative Commons Attribution 4.0 International licence, as recorded in the arXiv licence metadata for this exact version and shown on the abstract page https://arxiv.org/abs/2509.20836v3. Full rights notice: licenses/arxiv-2509.20836-CC-BY-4.0.txt.\par\smallskip

\noindent\textit{Editorial scope.} Counted unit: the Cor whose source label is \texttt{None} in the article (source0.tex lines 1063--1066, source label \texttt{None}), delivered together with the article's own complete original proof of it (source0.tex lines 1068--1107, one \texttt{proof} environment of 40 lines). No mathematics is written, repaired or re-derived: statement and proof are the article's own text line for line. Required context retained uncounted: prop:campbell (lines 374--378); prop:intensityli (lines 608--611); lem:separatedli (lines 625--627); eq:doubs (lines 813--815); prop:intmeasiden (lines 828--832); thm:finiteiv (lines 1057--1059). Every definition, display and label of those lines is retained unchanged. Omitted from this package and not redistributed: 1--373, 379--607, 612--624, 628--812, 816--827, 833--1056, 1060--1062, 1067--1067, 1108--1306. Nothing inside a retained excerpt is omitted or altered: every retained body is delivered exactly as the article's lines are written, so no omission receipt of any kind accompanies this package. Citations: the retained excerpts carry no citation command at all, so no bibliography entry of the article is transcribed and no reference is redistributed. Reference adaptation: none was needed and none was made; every reference inside the delivered unit resolves inside it, so no unresolved reference or citation is produced. Printed slips kept literal: no printed word, symbol, hypothesis or number of the counted statement or of its proof was altered. Layout and class: the article's own class is \texttt{amsart} with packages including geometry, amssymb, url, amsmath, amsfonts, amsthm, mathrsfs, appendix, inputenc, babel, cite, hyperref, cleveref, placeins; the wrapper is the lane's pinned \texttt{amsart} editorial wrapper at 10pt on A4, which restates the theorem-like environments the retained text needs and adds the article's own macro definitions that the retained text needs (\texttt{\textbackslash diag}), with extra wrapper packages (bm, bbm, dsfont, xspace, cancel, mathtools). The delivered numbering is the unit's own. Assets: no retained span contains a figure environment, a graphics-inclusion command, a PDF-page inclusion, a table or a TikZ picture; no image file is copied into this package, the package declares no asset dependency and no image file is redistributed with it. Licence: version arXiv:2509.20836v3 of this article is distributed under the Creative Commons Attribution 4.0 International licence, as recorded in the arXiv licence metadata for this exact version and shown on the abstract page https://arxiv.org/abs/2509.20836v3. A full rights notice is delivered at licenses/arxiv-2509.20836-CC-BY-4.0.txt. Characters: every retained excerpt is pure ASCII, so the delivered engine, XeLaTeX with its default Latin Modern text font, reports no missing character.
\begin{prop}[{\bf Campbell theorem}]\label{prop:campbell}
For every \(\mu\in\mathcal{M}^{G}\left(X\right)\), the transverse measure \(\mu_{Y}\) is \(\sigma\)-finite and satisfies
\[m_{G}\otimes\mu_{Y}\left(f\right)=\mu\left(f_{X}\right)\]
for every Borel function \(f:G\times Y\to\left[0,+\infty\right]\). In particular, \(\mu_{Y}\) is independent of the choice of \(w\).
\end{prop}

\begin{prop}\label{prop:intensityli}
Let \(\left(X,\mu\right)\) be a probability preserving \(G\)-space with a cross section \(Y\). Then \(Y\) is locally integrable precisely when the transverse measure \(\mu_{Y}\) is finite:
\[\iota_{\mu}\left(Y\right)\coloneqq \mu_{Y}\left(Y\right)<+\infty.\]
\end{prop}

\begin{lem}\label{lem:separatedli}
If \(Y\) is a separated cross section for a Borel \(G\)-space \(X\), then \(Y\) is locally integrable with respect to every \(G\)-invariant probability measures on \(X\).
\end{lem}

\begin{equation}\label{eq:doubs}\tag{\(\bigstar\)}
Y_{x_{1}}^{-1}Y_{x_{2}}\text{ is locally finite for all }x_{1},x_{2}\in X.
\end{equation}

\begin{prop}\label{prop:intmeasiden}
For every transverse \(G\)-space \(\left(X,\mu,Y\right)\) satisfying \eqref{eq:doubs}, the following identity holds:
\[\mu^{\left[r\right]}=\overline{\mu}_{Y^{\left[r\right]}}^{\left[r\right]};\]
that is, the intersection measure \(\mu^{\left[r\right]}\) is the transverse measure of the transverse \(G\)-space \(\big(\overline{Y}^{\left[r\right]},\overline{\mu}^{\left[r\right]},Y^{\left[r\right]}\big)\).
\end{prop}

\begin{thm}\label{thm:finiteiv}
Let \(\left(X,\mu,Y\right)\) be a transverse \(G\)-space and let \(r\in\mathbb{N}\). If \(e_{G}\in\Lambda_{Y}^{2r-1}\) is an isolated point, then \(I_{\mu}^{r}\left(Y\right)<+\infty\).
\end{thm}

\begin{cor}
Let \(\left(G,H;\Gamma\right)\) be a cut--and--project scheme, \(\left(\Xi,\xi\right)\) the associated probability preserving \(G\)-space, and \(Y_{W}\) the cross section for some window \(W\subset H\). Then
\[I_{\xi}^{r}\left(Y_{W}\right)<+\infty\text{ for all }r\in\mathbb{N}.\]
\end{cor}

\begin{proof}[Proof of Theorem~\ref{thm:finiteiv}]
Let \(\Lambda_{Y}\) be the return times set of \(Y\), and note that by the assumption \(\Lambda_{Y}^{2r-1}\), and a fortiori \(\Lambda_{Y}\), is locally finite. For the Borel \(G\)-space \(Y^{\left[r-1\right]}\) (in the diagonal action), define a new cross section \(\widetilde{Y}^{\left(r\right)}\) by
\[\widetilde{Y}^{\left(r\right)}\coloneqq\Lambda_{Y}^{r-1}\diag Y^{\otimes\left(r-1\right)}\subseteq Y^{\left[r-1\right]},\]
and note that its return times sets are
\[\widetilde{Y}_{\vec{z}}^{\left(r\right)}=\Lambda_{Y}^{r-1}Y_{\vec{z}}^{\otimes\left(r-1\right)}=\Lambda_{Y}^{r-1}\left(Y_{z_{2}}\cap\dotsm\cap Y_{z_{r}}\right)\text{ for }\vec{z}=\left(z_{2},\dotsc,z_{r}\right)\in Y^{\left[r-1\right]}.\]
Note also that for every \(\widetilde{y}\in\widetilde{Y}^{\left(r\right)}\), writing \(\widetilde{y}=h\diag\vec{y}\) for some \(h\in\Lambda_{Y}^{r-1}\) and \(\vec{y}\in Y^{\otimes\left(r-1\right)}\), we have
\[\widetilde{Y}_{\widetilde{y}}^{\left(r\right)}=\Lambda_{Y}^{r-1}Y_{\vec{y}}^{\otimes\left(r-1\right)}h^{-1}\subseteq\Lambda_{Y}^{2r-1},\]
since \(Y_{\vec{y}}^{\otimes\left(r-1\right)}\subseteq Y_{y_{1}}\subseteq\Lambda_{Y}\), and therefore
\[\Lambda_{\widetilde{Y}^{\left(r\right)}}\subseteq\Lambda_{Y}^{2r-1}.\]
From the assumption that \(e_{G}\in\Lambda_{Y}^{2r-1}\) is an isolated point it follows that \(\widetilde{Y}^{\left(r\right)}\) is a separated cross section, and therefore, using Lemma~\ref{lem:separatedli}, it is locally integrable with respect to \(\mu^{\left[r-1\right]}\), thus
\begin{equation}\label{eq:Ytilde}
\iota_{\mu^{\left[r-1\right]}}\big(\widetilde{Y}^{\left(r\right)}\big)=\mu_{\widetilde{Y}}^{\left[r-1\right]}\big(\widetilde{Y}^{\left(r\right)}\big)<+\infty.
\end{equation}

Recall the Borel \(G\)-space \(\overline{Y}^{\left[r\right]}\) with its cross section \(Y^{\left[r\right]}\). In order to relate its return times sets to those of \(\widetilde{Y}^{\left[r\right]}\), fix a Borel section \(\lambda:X\to G\) such that \(\lambda\left(x\right)\in Y_{x}\) for every \(x\in X\), and for every \(\left(x,\vec{z}\right)=\left(x,z_{2},\dotsc,z_{r}\right)\in\overline{Y}^{\left[r\right]}\) we find that
\begin{equation}\label{eq:Ytildereturn}
\begin{aligned}
Y_{\left(x,\vec{z}\right)}^{\left[r\right]}	
&=\left(Y_{z_{r}}^{-1}\cap\dotsm\cap Y_{z_{2}}^{-1}\right)Y_{x}=\big(Y_{\vec{z}}^{\otimes\left(r-1\right)}\big)^{-1}Y_{x}=\big(Y_{\vec{z}}^{\otimes\left(r-1\right)}\big)^{-1}Y_{\lambda\left(x\right).x}\lambda\left(x\right)\\
&\subseteq\big(Y_{\vec{z}}^{\otimes\left(r-1\right)}\big)^{-1}\Lambda_{Y}\lambda\left(x\right)=\big(\Lambda_{Y}Y_{\vec{z}}^{\otimes\left(r-1\right)}\big)^{-1}\lambda\left(x\right)=\big(\widetilde{Y}_{\vec{z}}^{\left(r\right)}\big)^{-1}\lambda\left(x\right).
\end{aligned}
\end{equation}
It follows that for every compact set \(K\subset G\),
\begin{align*}
&\int_{\overline{Y}^{\left[r\right]}}\big|Y_{\left(x,\vec{z}\right)}^{\left[r\right]}\cap K\big|d\overline{\mu}^{\left[r\right]}\left(x,\vec{z}\right)\\
\eqref{eq:Ytildereturn}&\qquad \leq\int_{\overline{Y}^{\left[r\right]}}\big|\big(\widetilde{Y}_{\vec{z}}^{\left(r\right)}\big)^{-1}\lambda\left(x\right)\cap K\big|d\overline{\mu}^{\left[r\right]}\left(x,\vec{z}\right)\\
&\qquad=\int_{\overline{Y}^{\left[r\right]}}\big|\widetilde{Y}_{\vec{z}}^{\left(r\right)}\cap\lambda\left(x\right)K^{-1}\big|d\overline{\mu}^{\left[r\right]}\left(x,\vec{z}\right)\\
&\qquad=\int_{X}\Big(\int_{Y^{\left[r-1\right]}}\big|\widetilde{Y}_{\vec{z}}^{\left(r\right)}\cap\lambda\left(x\right)K^{-1}\big|d\mu^{\left[r-1\right]}\left(\vec{z}\right)\Big)d\mu\left(x\right)\\
\left(\ast\right)&\qquad=\int_{X}\Big(m_{G}\big(K\lambda\left(x\right)^{-1}\big)\cdot\mu_{\widetilde{Y}}^{\left[r-1\right]}\big(\widetilde{Y}^{\left(r\right)}\big)\Big)d\mu\left(x\right)\\
\text{(unimodularity and }\eqref{eq:Ytilde})&\qquad=m_{G}\left(K\right)\cdot\mu_{\widetilde{Y}^{\left(r\right)}}^{\left[r-1\right]}\big(\widetilde{Y}^{\left(r\right)}\big)<+\infty,
\end{align*}
where \(\left(\ast\right)\) is by Campbell theorem~\ref{prop:campbell}, applied for each \(x\in X\) to the Borel function
\[f^{x}:G\times\widetilde{Y}^{\left(r\right)}\longrightarrow\left[0,+\infty\right],\quad f^{x}\left(g,\widetilde{y}\right)=\mathbf{1}_{K\lambda\left(x\right)^{-1}}\left(g\right),\]
whose \(Y^{\left[r-1\right]}\)-periodization (as in the proof of Proposition~\ref{prop:intensityli}) is
\[f_{Y^{\left[r-1\right]}}^{x}\left(\vec{z}\right)=\big|\widetilde{Y}_{\vec{z}}^{\left(r\right)}\cap\lambda\left(x\right)K^{-1}\big|.\]
This computation shows that \(Y^{\left[r\right]}\) is a locally integrable cross section for \(\big(\overline{Y}^{\left[r\right]},\overline{\mu}^{\left[r\right]}\big)\), which means that \(\overline{\mu}_{Y^{\left[r\right]}}^{\left[r\right]}\left(Y^{\left[r\right]}\right)<+\infty\) (recall Proposition~\ref{prop:intensityli}). Finally, in order to apply Proposition~\ref{prop:intmeasiden}, note that from \eqref{eq:Ytildereturn} for \(r=2\), for all \(\left(x,z\right)\in X^{\otimes 2}=\overline{Y}^{\left[2\right]}\) we have
\[Y_{x}^{-1}Y_{z}\subseteq\big(\widetilde{Y}_{z}^{\left(2\right)}\big)^{-1}\lambda\left(x\right),\]
and since \(\widetilde{Y}^{\left(2\right)}\) is locally integrable (\(e_{G}\in\Lambda_{Y}^{3}\subseteq\Lambda_{Y}^{2r-1}\) is an isolated point), \(Y_{x}^{-1}Y_{z}\) is a locally finite set, thus \eqref{eq:doubs} is satisfied. Then from Proposition~\ref{prop:intmeasiden} it follows that
\[I_{\mu}^{r}\left(Y\right)=\mu^{\left[r\right]}\big(Y^{\left[r\right]}\big)=\overline{\mu}_{Y^{\left[r\right]}}^{\left[r\right]}\big(Y^{\left[r\right]}\big)<+\infty.\qedhere\]
\end{proof}

\end{document}
